\section{Sintaxis y semántica del lenguaje}
\label{sec:sintaxis-semantica}

Hasta ahora las puertas se han introducido como notaciones para operadores concretos, y el signo $=$ se ha usado indistintamente para dos cosas distintas: que dos expresiones sean \emph{la misma expresión} y que \emph{representen el mismo operador}. Para poder razonar sobre reglas de reescritura (y demostrar que son correctas, que terminan o que son confluentes) es necesario separar ambos niveles. Esta sección lo hace en tres pasos: una \emph{sintaxis} que dice qué expresiones existen, una condición de \emph{buena formación} que descarta las que no tienen sentido, y una \emph{semántica denotacional} que asigna a cada expresión bien formada la matriz unitaria que representa.

A lo largo de la sección se fija el número de cúbits $n\in\N$. Como en el resto del trabajo, se sigue el convenio de Dirac: una expresión $\gate{C}\gate{D}$ aplicada a $\ket{a}$ se lee $\gate{C}(\gate{D}\ket{a})$, de modo que $\gate{D}$ actúa primero.

\subsection{Sintaxis abstracta}

La sintaxis se describe mediante una gramática en \emph{forma de Backus--Naur} (BNF, \emph{Backus--Naur Form}), la notación introducida en el informe de ALGOL~60 para gramáticas libres de contexto. Cada regla tiene la forma $\nt{símbolo} ::= \alpha_1 \mid \alpha_2 \mid \cdots$, que se lee «un $\nt{símbolo}$ es un $\alpha_1$, o un $\alpha_2$, o\ldots». Los símbolos $::=$ y $\mid$ pertenecen a la notación; las llaves, la coma, los paréntesis, $\emptyset$ y los nombres de las puertas son símbolos del propio lenguaje.

\begin{definition}[Sintaxis]
  \label{def:bnf}
  Los \define{términos} del lenguaje son las expresiones generadas por $\nt{circuito}$ en la gramática
  \[
    \begin{array}{rcl}
      \nt{circuito} & ::= & \gate{I} \mid \nt{puerta} \mid \nt{circuito}\,\nt{circuito} \mid (\,\nt{circuito}\,) \\[0.5ex]
      \nt{puerta}   & ::= & \nt{base}^{\,\nt{conjunto},\,\nt{conjunto}}_{\,\nt{conjunto}}                        \\[0.5ex]
      \nt{base}     & ::= & \gate{X} \mid \gate{Y} \mid \gate{Z} \mid \gate{H} \mid \gate{S} \mid \gate{T}        \\[0.5ex]
      \nt{conjunto} & ::= & \emptyset \mid \set{\,\nt{lista}\,}                                                 \\[0.5ex]
      \nt{lista}    & ::= & \nt{índice} \mid \nt{índice}\,,\,\nt{lista}                                         \\[0.5ex]
      \nt{índice}   & ::= & 1 \mid 2 \mid 3 \mid \cdots
    \end{array}
  \]
  Una puerta $\stategate{A}{\Gamma_1,\Gamma_0}{\Lambda}$ tiene \define{controles positivos} $\Gamma_1$, que la activan cuando el cúbit vale $\ket{1}$, \define{controles negativos} $\Gamma_0$, que la activan cuando vale $\ket{0}$, y \define{objetivos} $\Lambda$.
\end{definition}

Se trabaja con la \emph{sintaxis abstracta}: un término es el árbol de derivación de la gramática, no la cadena de símbolos que lo escribe. Los paréntesis solo sirven para escribir sin ambigüedad un árbol concreto, y por eso $(\gate{C}\gate{D})\gate{E}$ y $\gate{C}(\gate{D}\gate{E})$ son términos \emph{distintos}; que representen el mismo operador será una consecuencia de la semántica (proposición~\ref{prop:sem-asociatividad}), no de la sintaxis.

El conjunto de puertas base es una elección: la gramática se extiende sin cambios a cualquier conjunto finito $\mathcal{G}$ de puertas de un cúbit, o a familias parametrizadas como las rotaciones $R_z(\theta)$, añadiendo las alternativas correspondientes a $\nt{base}$.

\begin{notation}
  \label{not:abreviaturas}
  Se mantienen las abreviaturas de las secciones anteriores, que ahora son \emph{azúcar sintáctico} (formas cortas de escribir términos de la gramática):
  \begin{itemize}
    \item $\stategate{A}{\Gamma}{\Lambda}$ abrevia $\stategate{A}{\Gamma,\emptyset}{\Lambda}$, y $\gate{A}_{\Lambda}$ abrevia $\stategate{A}{\emptyset,\emptyset}{\Lambda}$;
    \item un índice $k$ abrevia el conjunto $\set{k}$, como en la sección de notación básica.
  \end{itemize}
  Así, por ejemplo, la puerta de Toffoli con controles en $1$ y $2$ y objetivo en $3$ es el término $\stategate{X}{1,2}{3}=\stategate{X}{\set{1,2},\emptyset}{\set{3}}$.
\end{notation}

\subsection{Buena formación}

La gramática genera expresiones como $\stategate{X}{\set{1},\emptyset}{\set{1}}$, en la que el cúbit $1$ controla y es objetivo a la vez, o $\gate{X}_{7}$ en un registro de tres cúbits. Ninguna de las dos representa una operación. Una gramática libre de contexto no puede expresar condiciones como «estos tres conjuntos son disjuntos» o «todo índice es como mucho $n$», así que se añaden como una segunda capa.

\begin{definition}[Buena formación]
  \label{def:buena-formacion}
  Se dice que un término $\gate{C}$ está \define{bien formado en $n$ cúbits}, y se escribe $\vdash_n \gate{C}$, cuando se deduce de las reglas siguientes:
  \begin{enumerate}
    \item $\vdash_n \gate{I}$;
    \item $\vdash_n \stategate{A}{\Gamma_1,\Gamma_0}{\Lambda}$ si $\Gamma_1,\Gamma_0,\Lambda\subseteq\N_n$ y son disjuntos dos a dos, es decir, $\Gamma_1\cap\Gamma_0=\Gamma_1\cap\Lambda=\Gamma_0\cap\Lambda=\emptyset$;
    \item $\vdash_n \gate{C}\gate{D}$ si $\vdash_n\gate{C}$ y $\vdash_n\gate{D}$.
  \end{enumerate}
\end{definition}

Se permite $\Lambda=\emptyset$; en tal caso, como se vio en la sección de puertas controladas, la puerta representa la identidad.

\subsection{Semántica denotacional}

La semántica asigna a cada término bien formado una matriz de $\C^{2^n\times 2^n}$. Se define \emph{por recursión sobre la gramática}: hay una cláusula por cada alternativa de la definición~\ref{def:bnf}, de modo que el significado de un término compuesto depende solo del significado de sus partes.

Para las puertas base se toman las matrices habituales:
\[
  \sem{\gate{X}}=\begin{pmatrix}0&1\\1&0\end{pmatrix},\quad
  \sem{\gate{Y}}=\begin{pmatrix}0&-i\\i&0\end{pmatrix},\quad
  \sem{\gate{Z}}=\begin{pmatrix}1&0\\0&-1\end{pmatrix},
\]
\[
  \sem{\gate{H}}=\frac{1}{\sqrt{2}}\begin{pmatrix}1&1\\1&-1\end{pmatrix},\quad
  \sem{\gate{S}}=\begin{pmatrix}1&0\\0&i\end{pmatrix},\quad
  \sem{\gate{T}}=\begin{pmatrix}1&0\\0&e^{i\pi/4}\end{pmatrix}.
\]

\begin{definition}
  \label{def:proyector-control}
  Dados $\Gamma_1,\Gamma_0\subseteq\N_n$ disjuntos, el \define{proyector de control} es
  \[
    \Pi_{\Gamma_1,\Gamma_0}=M_1\tensor\cdots\tensor M_n,
    \qquad
    M_j=\begin{cases}
      \dyad{1} & \text{si } j\in\Gamma_1, \\
      \dyad{0} & \text{si } j\in\Gamma_0, \\
      I        & \text{en otro caso.}
    \end{cases}
  \]
  Para una puerta base $\gate{A}$ y $\Lambda\subseteq\N_n$ se denota
  $\sem{\gate{A}}_\Lambda=\sem{\gate{A}}^{\delta_{1\Lambda}}\tensor\cdots\tensor\sem{\gate{A}}^{\delta_{n\Lambda}}$, es decir, $\sem{\gate{A}}$ en las posiciones de $\Lambda$ y la identidad en las demás.
\end{definition}

\begin{definition}[Semántica denotacional]
  \label{def:semantica}
  Para cada término $\gate{C}$ con $\vdash_n\gate{C}$ se define $\sem{\gate{C}}\in\C^{2^n\times 2^n}$ por:
  \begin{align*}
    \sem{\gate{I}}                                 & = I_{2^n},                                                                                   \\
    \sem{\stategate{A}{\Gamma_1,\Gamma_0}{\Lambda}} & = \left(I_{2^n}-\Pi_{\Gamma_1,\Gamma_0}\right)+\Pi_{\Gamma_1,\Gamma_0}\,\sem{\gate{A}}_\Lambda, \\
    \sem{\gate{C}\gate{D}}                         & = \sem{\gate{C}}\,\sem{\gate{D}}.
  \end{align*}
\end{definition}

La cláusula de las puertas dice exactamente lo esperado: sobre los estados en los que todos los controles toman el valor que los activa (la imagen de $\Pi_{\Gamma_1,\Gamma_0}$) se aplica $\gate{A}$ en los objetivos, y sobre el resto no se hace nada.

\begin{proposition}
  \label{prop:sem-unitaria}
  Si $\vdash_n\gate{C}$, entonces $\sem{\gate{C}}$ es una matriz unitaria.
\end{proposition}
\begin{proof}
  Por inducción sobre la estructura de $\gate{C}$. El caso $\gate{I}$ es inmediato, y el de $\gate{C}\gate{D}$ se sigue de que el producto de unitarias es unitario. Para una puerta, sean $\Pi=\Pi_{\Gamma_1,\Gamma_0}$ y $U=\sem{\gate{A}}_\Lambda$. Como $\Lambda$ es disjunto de $\Gamma_1\cup\Gamma_0$, $\Pi$ y $U$ actúan sobre factores tensoriales distintos y, por tanto, $\Pi U=U\Pi$ y $\Pi U^\dagger=U^\dagger\Pi$. Además $\Pi$ es un proyector ortogonal: $\Pi^\dagger=\Pi$, $\Pi^2=\Pi$ y $(I-\Pi)\Pi=0$. Entonces
  \begin{align*}
     & \left((I-\Pi)+\Pi U\right)^\dagger\left((I-\Pi)+\Pi U\right)                    \\
     & \qquad = (I-\Pi)^2 + (I-\Pi)\Pi U + U^\dagger\Pi(I-\Pi) + U^\dagger\Pi^2 U \\
     & \qquad = (I-\Pi) + \Pi\, U^\dagger U = (I-\Pi)+\Pi = I.
  \end{align*}
  En este cálculo es esencial la condición de buena formación: si $\Lambda$ y los controles se solaparan, $\Pi$ y $U$ no conmutarían en general y la expresión de la definición~\ref{def:semantica} podría no ser unitaria.
\end{proof}

\begin{proposition}
  \label{prop:sem-coherente}
  La semántica es coherente con la definición de puerta controlada de la sección de puertas controladas. Si $\Gamma=\set{h_1,\dots,h_l}$ y $\ket{a}=\ket{a_1\cdots a_n}$ es un ket producto, entonces
  \[
    \sem{\stategate{A}{\Gamma}{\Lambda}}\ket{a}
    = \ket{a}-a_{h_{1}1}\cdots a_{h_{l}1}\stateket{a}{\Gamma}{1}+a_{h_{1}1}\cdots a_{h_{l}1}\,\sem{\gate{A}}_\Lambda\stateket{a}{\Gamma}{1}.
  \]
\end{proposition}
\begin{proof}
  Por la definición~\ref{def:proyector-control}, $\Pi_{\Gamma,\emptyset}$ actúa factor a factor: en cada $h\in\Gamma$ lleva $\ket{a_h}=a_{h0}\ket{0}+a_{h1}\ket{1}$ a $a_{h1}\ket{1}$, y deja los demás factores igual. Por tanto $\Pi_{\Gamma,\emptyset}\ket{a}=a_{h_{1}1}\cdots a_{h_{l}1}\stateket{a}{\Gamma}{1}$, y basta sustituir en la definición~\ref{def:semantica}.
\end{proof}

\begin{example}
  Para $n=3$, el término $\stategate{X}{1,2}{3}$ está bien formado y
  $\sem{\stategate{X}{1,2}{3}}=(I-\dyad{1}\tensor\dyad{1}\tensor I)+\dyad{1}\tensor\dyad{1}\tensor\sem{\gate{X}}$,
  que es la matriz de la puerta de Toffoli: intercambia $\ket{110}$ y $\ket{111}$ y deja fijos los demás vectores de la base. En cambio, $\stategate{X}{1,\emptyset}{1}$ no está bien formado y no tiene semántica.
\end{example}

\subsection{Igualdad sintáctica e igualdad semántica}

\begin{definition}
  \label{def:equivalencia}
  Dos términos $\gate{C},\gate{D}$ con $\vdash_n\gate{C}$ y $\vdash_n\gate{D}$ son \define{semánticamente equivalentes}, y se escribe $\gate{C}\cong\gate{D}$, si $\sem{\gate{C}}=\sem{\gate{D}}$.
\end{definition}

Así quedan separados dos niveles: dos términos pueden ser distintos como árboles y, sin embargo, equivalentes. Las igualdades entre puertas demostradas en las secciones anteriores, como~\eqref{eq:conmutacion-gates} o~\eqref{eq:puerta-involutiva}, son en este lenguaje equivalencias $\cong$ entre términos distintos.

\begin{proposition}
  \label{prop:sem-asociatividad}
  Para términos bien formados, $(\gate{C}\gate{D})\gate{E}\cong\gate{C}(\gate{D}\gate{E})$, $\gate{I}\gate{C}\cong\gate{C}$ y $\gate{C}\gate{I}\cong\gate{C}$.
\end{proposition}
\begin{proof}
  Por la definición~\ref{def:semantica}, se reduce a la asociatividad del producto de matrices y a que $I_{2^n}$ es su neutro.
\end{proof}

En lo sucesivo, por la proposición~\ref{prop:sem-asociatividad}, se omiten los paréntesis en las composiciones.

Una \define{regla de reescritura} es un par de términos $\gate{L}\to\gate{R}$, posiblemente con conjuntos de índices como parámetros, junto con las condiciones sobre esos parámetros bajo las que se aplica. La regla es \define{correcta} (\emph{sound}) si $\gate{L}\cong\gate{R}$ para todos los valores de los parámetros que cumplen esas condiciones y hacen que ambos lados estén bien formados. La corrección de las reglas garantiza que reescribir un circuito nunca cambia el operador que representa.

Como ejemplo, la regla siguiente cambia un control positivo por uno negativo conjugando con $\gate{X}$, y es la base para relacionar las puertas con controles negativos con las de la sección de puertas controladas.

\begin{proposition}[Inversión de un control]
  \label{prop:inversion-control}
  Sean $\Gamma_1,\Gamma_0,\Lambda\subseteq\N_n$ disjuntos dos a dos y $c\in\N_n\setminus(\Gamma_1\cup\Gamma_0\cup\Lambda)$. Entonces
  \[
    \gate{X}_c\,\stategate{A}{\Gamma_1\cup\set{c},\,\Gamma_0}{\Lambda}\,\gate{X}_c
    \;\cong\;
    \stategate{A}{\Gamma_1,\,\Gamma_0\cup\set{c}}{\Lambda}.
  \]
\end{proposition}
\begin{proof}
  Sean $P=\sem{\gate{X}_c}$, $\Pi=\Pi_{\Gamma_1\cup\set{c},\Gamma_0}$ y $U=\sem{\gate{A}}_\Lambda$. Como $c\notin\Lambda$, $P$ y $U$ actúan sobre factores tensoriales distintos y conmutan; además $P^2=I$. Por tanto
  \[
    P\left((I-\Pi)+\Pi U\right)P = (I-P\Pi P) + P\Pi P\, U.
  \]
  El proyector $P\Pi P$ coincide con $\Pi$ en todos los factores salvo el $c$-ésimo, donde cambia $\dyad{1}$ por $\sem{\gate{X}}\dyad{1}\sem{\gate{X}}=\dyad{0}$. Es decir, $P\Pi P=\Pi_{\Gamma_1,\Gamma_0\cup\set{c}}$, y el lado derecho es $\sem{\stategate{A}{\Gamma_1,\Gamma_0\cup\set{c}}{\Lambda}}$.
\end{proof}

\begin{corollary}
  \label{cor:controles-negativos}
  En las condiciones anteriores,
  $\stategate{A}{\Gamma_1,\Gamma_0}{\Lambda}\cong\gate{X}_{\Gamma_0}\,\stategate{A}{\Gamma_1\cup\Gamma_0}{\Lambda}\,\gate{X}_{\Gamma_0}$.
  En particular, toda puerta con controles negativos es equivalente a un término que solo usa controles positivos.
\end{corollary}
\begin{proof}
  Por inducción sobre $|\Gamma_0|$, aplicando la proposición~\ref{prop:inversion-control} a cada $c\in\Gamma_0$ y usando que las puertas $\gate{X}_c$ con $c\in\Gamma_0$ conmutan entre sí~\eqref{eq:conmutacion-gates}.
\end{proof}

Con esta base, los pasos siguientes de la formalización son fijar un conjunto de reglas orientadas, demostrar que cada una es correcta en el sentido anterior y estudiar si el sistema resultante termina y es confluente, lo que daría una forma normal para cada circuito. La pregunta inversa, si las reglas bastan para deducir \emph{toda} equivalencia $\gate{C}\cong\gate{D}$ (la \emph{completitud}), es la que resuelven para los circuitos cuánticos generales \textcite{Clement2023} y, para los circuitos Toffoli, \textcite{CockettComfort2018}.
